\documentclass[]{article}

\usepackage{geometry}
\usepackage{longtable}
\usepackage{listings}
\usepackage{indentfirst}
\usepackage{multirow}

\geometry{
	a4paper,
	total={170mm,257mm},
	left=20mm,
	top=20mm,
}


%opening
\title{ \textbf{Dancing Pipes} \vspace{2em} \\ Programmierprojekt \\ Sommersemester 2024}


\author{ Atilla Özdemir \and Kirill Dedov \and Matsvei Tsiareshka \and Toni Kostov}

\begin{document}

\maketitle
 \thispagestyle{empty}
\newpage

\setcounter{page}{1}

\section{So macht man eine Section}
Hier kann man text schreiben.

So macht man eine neue Zeile.
So kann man Code im Text schreiben: zum Beispiel \texttt{int a = 5;}.

\subsection{So macht man eine Subsection}

Wenn man mehrere Zeilen Code schreiben will nutzt man das:

\begin{lstlisting}[language=Java]
	public static void changeChord() {
		sequencer.stop();
		sequencer.setSequence(sequence);
		sequencer.start();
	}
\end{lstlisting}

\subsection*{So macht man eine Subsection ohne abzählen}
Tabellen kann man auch machen:
\begin{center}
	\begin{longtable}{|l|l|}
		\textbf{Spalte1} & 
		\textbf{Spalte 2} \\
		\hline
		
		hier kann man text schreiben & und hier auch \\
		\hline
		
		noch Text schreiben & 
		und hier auch mehr text  \\
		\hline
		
		\multicolumn{2}{|l|}{Spalten kann man auch so mergen} \\
		\hline
		
		und hier & weiter machen \\
		\hline
		\hline
		
		\multirow[c]{1}{*}{Zeilen kann man auch mergen} & das geht so  \\
	
		& und dann so \\ 
		\hline
		
		danach & einfach weiter machen \\
		\hline
		
	\end{longtable}
\end{center}

Wenn man eine Liste machen will, macht man das so:

\begin{itemize}
	\item Element 1
	\item Element 2
	\item Element 3
	\item Element 4
\end{itemize}

Viel Spaß :)



\subsubsection{Man kann sogar Subsubsection machen}

	
\section{Kamera}

\section{MIDI}

\section{Server}

\end{document}
